% Software Overview
% MIT License
%
% Copyright (c) 2026 HumberASV
% Copyright (c) 2026 Carson Fujita
%
% Permission is hereby granted, free of charge, to any person obtaining a copy
% of this software and associated documentation files (the "Software"), to deal
% in the Software without restriction, including without limitation the rights
% to use, copy, modify, merge, publish, distribute, sublicense, and/or sell
% copies of the Software, and to permit persons to whom the Software is
% furnished to do so, subject to the following conditions:
%
% The above copyright notice and this permission notice shall be included in all
% copies or substantial portions of the Software.
%
% THE SOFTWARE IS PROVIDED "AS IS", WITHOUT WARRANTY OF ANY KIND, EXPRESS OR
% IMPLIED, INCLUDING BUT NOT LIMITED TO THE WARRANTIES OF MERCHANTABILITY,
% FITNESS FOR A PARTICULAR PURPOSE AND NONINFRINGEMENT. IN NO EVENT SHALL THE
% AUTHORS OR COPYRIGHT HOLDERS BE LIABLE FOR ANY CLAIM, DAMAGES OR OTHER
% LIABILITY, WHETHER IN AN ACTION OF CONTRACT, TORT OR OTHERWISE, ARISING FROM,
% OUT OF OR IN CONNECTION WITH THE SOFTWARE OR THE USE OR OTHER DEALINGS IN THE
% SOFTWARE.

% This remote connections for Loon-E ASV
% This is separated from the development section as this may change more frequently.

\subsection{Primary Computer}
\label{sec:primary-computer}

The primary computer is the main processing unit of the ASV, responsible for controlling the vehicle and processing data from its sensors. The software running on the primary computer is developed using ROS 2, a popular robotics middleware that provides a framework for building robot applications \autocite{ROSWIKI}.

We use the Jetson Orin Nano as the primary computer for Loon-E.

Details on development for the primary computer can be found in section \ref{sec:active-development} which covers the active development process for the software team, including how to set up the development environment, how to run the software, and how to contribute to the project. It also includes any relevant documentation or resources for developers working on the primary computer software.